%%% =======================================================================
%%% ISLS Extension Phase 4 v1.0.0
%%% C19: Universal Gateway
%%% C20: Domain Adapter Library
%%% Author: Sebastian Klemm
%%% Date: 16 March 2026
%%% Status: Implementation Specification for Claude Code
%%% =======================================================================
\documentclass[11pt,a4paper,twoside]{article}

\usepackage[utf8]{inputenc}
\usepackage{amsmath,amssymb,amsthm}
\usepackage{algorithm,algorithmic}
\usepackage{booktabs,longtable,tabularx}
\usepackage{enumitem}
\usepackage{geometry}
\usepackage{hyperref}
\usepackage{cleveref}
\usepackage{xcolor}
\usepackage{fancyhdr}
\usepackage{listings}

\geometry{margin=2.5cm,top=3cm,bottom=3cm,headheight=14pt}

\theoremstyle{definition}
\newtheorem{definition}{Definition}[section]
\newtheorem{requirement}{Requirement}[section]
\newtheorem{invariant}{Invariant}[section]

\theoremstyle{plain}
\newtheorem{theorem}{Theorem}[section]
\newtheorem{proposition}{Proposition}[section]

\theoremstyle{remark}
\newtheorem{remark}{Remark}[section]

\newcommand{\ISLS}{\textsf{ISLS}}
\newcommand{\RR}{\mathbb{R}}
\DeclareMathOperator{\Dig}{Dig}
\DeclareMathOperator{\Canon}{Canon}

\lstset{
  basicstyle=\ttfamily\small,
  breaklines=true,
  frame=single,
  backgroundcolor=\color{gray!5},
  rulecolor=\color{gray!30},
}

\pagestyle{fancy}
\fancyhf{}
\fancyhead[LE,RO]{\ISLS{} Extension Phase 4 v1.0.0}
\fancyhead[RE,LO]{Sebastian Klemm}
\fancyfoot[C]{\thepage}

\hypersetup{
  colorlinks=true,
  linkcolor=blue!60!black,
  pdftitle={ISLS Extension Phase 4 v1.0.0},
  pdfauthor={Sebastian Klemm},
}

\begin{document}

\begin{titlepage}
\centering
\vspace*{2.5cm}
{\Huge\bfseries \ISLS{} Extension Phase~4\\[0.3cm]
v1.0.0}\\[1.5cm]
{\Large C19: Universal Gateway\\
C20: Domain Adapter Library}\\[1cm]
{\large The Plug-and-Play Interface for the\\
Intelligent Semantic Ledger Substrate}\\[2cm]
{\large Sebastian Klemm}\\[0.5cm]
{\large 16 March 2026}\\[2cm]

\begin{abstract}
\noindent
This document specifies the external interface layer for \ISLS{}: a
Universal Gateway (C19) providing a single unified API for data
ingestion, crystal querying, and live streaming; and a Domain Adapter
Library (C20) providing ready-made connectors for common data sources
and sinks. Together, these two crates transform \ISLS{} from a
command-line tool into an embeddable, network-accessible service that
any system can connect to without understanding \ISLS{} internals.

C19 absorbs the Node API pattern from OMNIPOL/ABRAXAS and the
Filament adapter trait. C20 absorbs the Ingestion Matrix concept from
MEF-Core and the Domain Adapter pattern from SPEC-003.

After Phase~4, \ISLS{} is operationally complete: data flows in from
arbitrary sources, crystals flow out to arbitrary consumers, and the
entire system is controllable via a single REST/WebSocket endpoint.

\medskip\noindent\textbf{Status:} Implementation Specification.\\
\textbf{Parent:} ISLS v1.0.0 + Extensions Phase 1--3.\\
\textbf{Absorbs:} OMNIPOL Node API, MEF-Core Ingestion Matrix,
  SPEC-003 Domain Adapter.
\end{abstract}
\end{titlepage}

\tableofcontents
\newpage


%%% =====================================================================
\part{C19: Universal Gateway}\label{part:gateway}

\section{Purpose}\label{sec:gw:purpose}

The gateway is the single external interface to a running \ISLS{}
instance. It handles three flows:

\begin{enumerate}[nosep,label=(\alph*)]
  \item \textbf{Ingest}: accept observations from external sources,
        canonicalize them, and feed them into the engine.
  \item \textbf{Query}: expose crystals, manifests, metrics, alerts,
        and validation results to external consumers.
  \item \textbf{Stream}: provide real-time event feeds (crystal
        emissions, gate decisions, metric updates, scale events) via
        WebSocket or Server-Sent Events.
\end{enumerate}

One gateway. One port. Three flows. Everything else is an adapter
(C20) that translates domain-specific protocols into gateway calls.


\section{Architecture}\label{sec:gw:arch}

\begin{lstlisting}[title={Gateway Position in the Stack}]
External World
    |
    v
[isls-gateway]  <-- C19: REST + WebSocket server
    |
    +---> [isls-adapters]  <-- C20: domain connectors
    |
    +---> [isls-engine]    <-- observation ingestion
    |
    +---> [isls-store]     <-- crystal/manifest queries
    |
    +---> [isls-archive]   <-- crystal retrieval
    |
    +---> [isls-manifest]  <-- manifest retrieval
    |
    +---> [isls-capsule]   <-- seal/open operations
\end{lstlisting}

\begin{requirement}[Single Binary]\label{req:singlebin}
The gateway \textbf{MUST} be embeddable in the existing \texttt{isls}
binary via a \texttt{serve} subcommand:
\begin{lstlisting}
isls serve --port 8420 --host 0.0.0.0
\end{lstlisting}
No separate binary, no separate process. One binary, one command.
\end{requirement}

\begin{requirement}[Stateless Gateway]\label{req:stateless}
The gateway itself \textbf{MUST NOT} maintain state beyond request
handling. All state lives in the engine, store, and archive. The
gateway is a translation layer, not a service with its own database.
\end{requirement}


\section{REST API}\label{sec:gw:rest}

\subsection{System Endpoints}

\begin{longtable}{l l p{6.5cm}}
\toprule
\textbf{Method} & \textbf{Path} & \textbf{Purpose} \\
\midrule
\endfirsthead
\midrule
\endhead
\bottomrule
\endfoot
\texttt{GET} & \texttt{/health} & Liveness check. Returns version,
  uptime, engine state, crystal count, entity count. \\
\texttt{GET} & \texttt{/config} & Current RunDescriptor and
  configuration (read-only). \\
\texttt{GET} & \texttt{/status} & Engine status: running/idle, current
  tick, active adapters, scale state. \\
\end{longtable}

\subsection{Ingestion Endpoints}

\begin{longtable}{l l p{6.5cm}}
\toprule
\textbf{Method} & \textbf{Path} & \textbf{Purpose} \\
\midrule
\endfirsthead
\midrule
\endhead
\bottomrule
\endfoot
\texttt{POST} & \texttt{/ingest} & Submit a batch of observations.
  Body: JSON array of observation objects. Returns: ingestion receipt
  with observation count, duplicate count, digest. \\
\texttt{POST} & \texttt{/ingest/stream} & Open a streaming ingestion
  channel. Body: newline-delimited JSON (NDJSON). Each line is one
  observation. Returns: per-line receipts. \\
\texttt{POST} & \texttt{/ingest/file} & Upload a file for ingestion.
  Multipart form with file field. Adapter auto-detected from
  Content-Type or file extension. \\
\texttt{GET} & \texttt{/ingest/adapters} & List available ingestion
  adapters with their status and configuration. \\
\end{longtable}

\subsection{Query Endpoints}

\begin{longtable}{l l p{6.5cm}}
\toprule
\textbf{Method} & \textbf{Path} & \textbf{Purpose} \\
\midrule
\endfirsthead
\midrule
\endhead
\bottomrule
\endfoot
\texttt{GET} & \texttt{/crystals} & List crystals. Query params:
  \texttt{?scale=micro\&limit=50\&after=<tick>}. Returns: array of
  crystal summaries. \\
\texttt{GET} & \texttt{/crystals/\{id\}} & Get full crystal by ID.
  Returns: complete crystal JSON including topology signature,
  evidence chain, constraint program. \\
\texttt{GET} & \texttt{/crystals/\{id\}/validate} & Run formal
  validation on a specific crystal. Returns: 8-check breakdown. \\
\texttt{GET} & \texttt{/crystals/\{id\}/replay} & Download replay pack
  for this crystal's parent run. \\
\texttt{GET} & \texttt{/manifests} & List execution manifests.
  Query params: \texttt{?project=<name>\&limit=10}. \\
\texttt{GET} & \texttt{/manifests/\{run\_id\}} & Get full manifest by
  run ID. \\
\texttt{GET} & \texttt{/manifests/\{run\_id\}/verify} & Verify
  manifest integrity (MV1--MV6). \\
\texttt{GET} & \texttt{/metrics} & Get latest metric snapshot.
  Query params: \texttt{?run=<id>\&tick=<n>}. \\
\texttt{GET} & \texttt{/metrics/history} & Get metric time series.
  Query params: \texttt{?metric=M25\&from=<tick>\&to=<tick>}. \\
\texttt{GET} & \texttt{/alerts} & List active alerts. Query params:
  \texttt{?level=red\&limit=20}. \\
\texttt{GET} & \texttt{/scale} & Get current multi-scale state:
  cluster count, bridge activity, domain count. \\
\texttt{GET} & \texttt{/scale/universes} & List all Hypercube
  Universes with aggregate states. \\
\texttt{GET} & \texttt{/scale/bridges} & List all bridges with
  activation status. \\
\texttt{GET} & \texttt{/topology} & Get current spectral
  decomposition summary: spectral gap, Kuramoto $r$, Cheeger
  estimate. \\
\end{longtable}

\subsection{Action Endpoints}

\begin{longtable}{l l p{6.5cm}}
\toprule
\textbf{Method} & \textbf{Path} & \textbf{Purpose} \\
\midrule
\endfirsthead
\midrule
\endhead
\bottomrule
\endfoot
\texttt{POST} & \texttt{/engine/start} & Start the engine loop.
  Body: \texttt{\{mode, ticks, project\}}. \\
\texttt{POST} & \texttt{/engine/stop} & Gracefully stop the engine
  after current tick completes. \\
\texttt{POST} & \texttt{/engine/step} & Execute exactly one
  macro-step. Returns: tick result. \\
\texttt{POST} & \texttt{/execute} & Execute a crystal or program in
  generative mode. Body: crystal JSON or HDAG. \\
\texttt{POST} & \texttt{/capsule/seal} & Seal a secret. Body:
  \texttt{\{secret, lock, master\_key\_ref\}}. \\
\texttt{POST} & \texttt{/capsule/open} & Open a capsule. Body:
  capsule JSON + master key reference. \\
\texttt{POST} & \texttt{/validate/formal} & Run formal validation on
  all crystals in archive. Returns: summary. \\
\texttt{GET} & \texttt{/report/html} & Generate and return the full
  HTML validation report. \\
\texttt{GET} & \texttt{/export/\{run\_id\}} & Export a complete run as
  ZIP (manifest + crystals + traces + metrics). \\
\end{longtable}

\subsection{Response Format}

\begin{definition}[Envelope]\label{def:envelope}
All responses use a standard envelope:
\begin{lstlisting}
{
  "ok": true,
  "data": { ... },
  "meta": {
    "request_id": "...",
    "timestamp": "2026-03-16T...",
    "duration_ms": 42
  }
}
\end{lstlisting}
Error responses:
\begin{lstlisting}
{
  "ok": false,
  "error": {
    "code": "CRYSTAL_NOT_FOUND",
    "message": "No crystal with ID abc123..."
  },
  "meta": { ... }
}
\end{lstlisting}
\end{definition}

\begin{definition}[Error Codes]\label{def:errors}
Standard error codes:
\begin{itemize}[nosep]
  \item \texttt{400 BAD\_REQUEST}: malformed input.
  \item \texttt{404 NOT\_FOUND}: crystal, manifest, or run not found.
  \item \texttt{409 CONFLICT}: determinism conflict (parity failure).
  \item \texttt{422 VALIDATION\_FAILED}: observation rejected by
        canonicalization or deduplication.
  \item \texttt{429 RATE\_LIMITED}: ingestion rate exceeded.
  \item \texttt{499 GATE\_ABORT}: engine stopped due to gate failure
        or budget exhaustion.
  \item \texttt{503 ENGINE\_BUSY}: engine is already running, cannot
        start a second run.
\end{itemize}
\end{definition}


\section{WebSocket API}\label{sec:gw:ws}

\begin{definition}[Event Stream]\label{def:eventstream}
A WebSocket connection to \texttt{ws://host:port/events} delivers a
real-time event feed. Each message is a JSON object with a
\texttt{type} field:

\begin{lstlisting}
{ "type": "crystal",     "data": { crystal summary } }
{ "type": "gate",        "data": { gate_snapshot } }
{ "type": "metric",      "data": { metric_snapshot } }
{ "type": "alert",       "data": { alert } }
{ "type": "scale_event", "data": { cluster_split/merge/bridge } }
{ "type": "fixpoint",    "data": { scale, iteration } }
{ "type": "tick",        "data": { tick_number, entities, crystals } }
\end{lstlisting}
\end{definition}

\begin{definition}[Subscription Filter]\label{def:subfilter}
Clients \textbf{MAY} send a subscription message to filter events:
\begin{lstlisting}
{ "subscribe": ["crystal", "alert", "metric"] }
\end{lstlisting}
Default (no subscription message): all event types.
\end{definition}

\begin{requirement}[Backpressure]\label{req:backpressure}
If a WebSocket client cannot keep up, the gateway \textbf{MUST}
buffer up to 1000 events and then drop oldest events. Dropped events
are counted in the \texttt{/health} endpoint as
\texttt{ws\_dropped\_events}.
\end{requirement}


\section{Authentication}\label{sec:gw:auth}

\begin{requirement}[Optional API Key]\label{req:apikey}
Authentication is optional and disabled by default (single-operator
use case). When enabled via configuration, the gateway \textbf{MUST}
require a bearer token in the \texttt{Authorization} header:
\begin{lstlisting}
Authorization: Bearer <api_key>
\end{lstlisting}
The API key is stored as a SHA-256 hash in the configuration. No
user/role management. One key, one operator.
\end{requirement}

\begin{requirement}[CORS]\label{req:cors}
The gateway \textbf{MUST} support configurable CORS headers for
browser-based clients:
\begin{lstlisting}
gateway:
  cors:
    allowed_origins: ["http://localhost:*"]
    allowed_methods: ["GET", "POST"]
\end{lstlisting}
\end{requirement}


\section{Configuration}\label{sec:gw:config}

\begin{lstlisting}[title={Gateway Configuration (in isls config.yaml)}]
gateway:
  enabled: true
  host: "127.0.0.1"
  port: 8420
  max_connections: 100
  request_timeout_ms: 30000
  ingestion_rate_limit: 10000  # observations per second
  auth:
    enabled: false
    api_key_hash: ""  # SHA-256 of the API key
  cors:
    enabled: false
    allowed_origins: []
  ws:
    max_buffer: 1000
    heartbeat_interval_ms: 30000
\end{lstlisting}


\section{Rust API}\label{sec:gw:rustapi}

\begin{lstlisting}[title={isls-gateway/src/lib.rs}]
pub struct GatewayConfig {
    pub host: String,
    pub port: u16,
    pub max_connections: usize,
    pub request_timeout_ms: u64,
    pub ingestion_rate_limit: u64,
    pub auth: AuthConfig,
    pub cors: CorsConfig,
    pub ws: WsConfig,
}

pub struct Gateway {
    config: GatewayConfig,
    engine: Arc<Mutex<Engine>>,
    store: Arc<IslandStore>,
    archive: Arc<Archive>,
    adapters: Arc<AdapterRegistry>,
}

impl Gateway {
    pub fn new(
        config: GatewayConfig,
        engine: Engine,
        store: IslandStore,
        archive: Archive,
        adapters: AdapterRegistry,
    ) -> Self;

    /// Start the HTTP/WS server (blocking)
    pub async fn serve(&self) -> Result<()>;

    /// Graceful shutdown
    pub async fn shutdown(&self) -> Result<()>;
}

// Ingestion receipt
pub struct IngestReceipt {
    pub observations_received: u64,
    pub observations_accepted: u64,
    pub duplicates_skipped: u64,
    pub digest: Hash256,
    pub timestamp: f64,
}

// Event types for WebSocket
pub enum GatewayEvent {
    Crystal(CrystalSummary),
    Gate(GateSnapshot),
    Metric(MetricSnapshot),
    Alert(AlertRow),
    ScaleEvent(ScaleEvent),
    Fixpoint { scale: Scale, iteration: u64 },
    Tick { number: u64, entities: u64, crystals: u64 },
}
\end{lstlisting}


%%% =====================================================================
\part{C20: Domain Adapter Library}\label{part:adapters}

\section{Purpose}\label{sec:ad:purpose}

Adapters translate between the external world and \ISLS{}'s canonical
observation format. Each adapter implements a single trait and handles
one data source or sink. The adapter library ships with built-in
adapters for common use cases so that users can connect \ISLS{} to
their data without writing code.


\section{The Filament Trait}\label{sec:ad:filament}

Absorbed from OMNIPOL's Filament concept and MEF-Core's Ingestion
Matrix.

\begin{definition}[Ingestion Adapter Trait]\label{def:ingesttrait}
Every ingestion adapter implements:
\begin{lstlisting}
pub trait IngestionAdapter: Send + Sync {
    /// Human-readable name
    fn name(&self) -> &str;

    /// Which content types / file extensions this adapter handles
    fn accepts(&self) -> Vec<String>;

    /// Transform raw bytes into canonical observations
    fn transform(
        &self,
        input: &[u8],
        config: &AdapterConfig,
    ) -> Result<Vec<Observation>>;

    /// For streaming adapters: start a background source
    fn start_stream(
        &self,
        config: &AdapterConfig,
        sender: Sender<Vec<Observation>>,
    ) -> Result<JoinHandle<()>> {
        // Default: not a streaming adapter
        Err(AdapterError::NotStreaming)
    }

    /// Healthcheck
    fn health(&self) -> AdapterHealth;
}
\end{lstlisting}
\end{definition}

\begin{definition}[Emission Adapter Trait]\label{def:emittrait}
Every emission adapter (output sink) implements:
\begin{lstlisting}
pub trait EmissionAdapter: Send + Sync {
    /// Human-readable name
    fn name(&self) -> &str;

    /// Receive a crystal event and forward it
    fn emit_crystal(
        &self,
        crystal: &SemanticCrystal,
        config: &AdapterConfig,
    ) -> Result<()>;

    /// Receive an alert and forward it
    fn emit_alert(
        &self,
        alert: &AlertRow,
        config: &AdapterConfig,
    ) -> Result<()>;

    /// Receive a metric snapshot and forward it
    fn emit_metric(
        &self,
        snapshot: &MetricSnapshot,
        config: &AdapterConfig,
    ) -> Result<()>;
}
\end{lstlisting}
\end{definition}

\begin{definition}[Adapter Registry]\label{def:adapterreg}
The adapter registry holds all available adapters:
\begin{lstlisting}
pub struct AdapterRegistry {
    ingestion: BTreeMap<String, Box<dyn IngestionAdapter>>,
    emission: BTreeMap<String, Box<dyn EmissionAdapter>>,
}

impl AdapterRegistry {
    pub fn new() -> Self;
    pub fn register_ingestion(&mut self, a: Box<dyn IngestionAdapter>);
    pub fn register_emission(&mut self, a: Box<dyn EmissionAdapter>);
    pub fn get_ingestion(&self, name: &str) -> Option<&dyn IngestionAdapter>;
    pub fn auto_detect(&self, content_type: &str) -> Option<&dyn IngestionAdapter>;
    pub fn list_ingestion(&self) -> Vec<&str>;
    pub fn list_emission(&self) -> Vec<&str>;

    /// Register all built-in adapters
    pub fn register_defaults(&mut self);
}
\end{lstlisting}
\end{definition}

\begin{invariant}[Adapter Determinism]\label{inv:adapterdet}
Every ingestion adapter's \texttt{transform} method \textbf{MUST} be
deterministic: identical input bytes with identical config produce
identical observations. Adapters \textbf{MUST NOT} use
\texttt{rand}, system time, or network calls inside
\texttt{transform}. Timestamps in observations come from the data
itself or from the gateway's ingestion timestamp, not from the
adapter.
\end{invariant}


\section{Built-In Ingestion Adapters}\label{sec:ad:ingestion}

\subsection{A01: JSON Adapter}

\begin{tabularx}{\textwidth}{l X}
\toprule
\textbf{Name} & \texttt{json} \\
\textbf{Accepts} & \texttt{application/json}, \texttt{.json}, \texttt{.jsonl} \\
\textbf{Input} & JSON array of objects, or newline-delimited JSON \\
\textbf{Transform} & Each JSON object becomes one observation. Entity
  ID from \texttt{id} or \texttt{entity\_id} field. Payload is the
  full object (canonical JSON). \\
\textbf{Config} & \texttt{entity\_id\_field}: which field to use as
  entity ID (default: \texttt{"id"}). \\
\bottomrule
\end{tabularx}

\subsection{A02: CSV Adapter}

\begin{tabularx}{\textwidth}{l X}
\toprule
\textbf{Name} & \texttt{csv} \\
\textbf{Accepts} & \texttt{text/csv}, \texttt{.csv}, \texttt{.tsv} \\
\textbf{Input} & CSV with header row \\
\textbf{Transform} & Each row becomes one observation. Entity ID from
  configurable column. Numeric columns parsed as f64. \\
\textbf{Config} & \texttt{entity\_id\_column}, \texttt{delimiter},
  \texttt{timestamp\_column} (optional). \\
\bottomrule
\end{tabularx}

\subsection{A03: Syslog / Structured Log Adapter}

\begin{tabularx}{\textwidth}{l X}
\toprule
\textbf{Name} & \texttt{syslog} \\
\textbf{Accepts} & \texttt{text/plain} (syslog format),
  \texttt{application/x-ndjson} (structured logs) \\
\textbf{Input} & Syslog lines (RFC 5424) or structured JSON logs \\
\textbf{Transform} & Entity ID from hostname + process name. Payload
  includes severity, facility, message hash. Structured logs use
  their own field names. \\
\textbf{Config} & \texttt{format}: \texttt{rfc5424} or
  \texttt{structured}. \\
\bottomrule
\end{tabularx}

\subsection{A04: Cryptocurrency Market Data Adapter}

\begin{tabularx}{\textwidth}{l X}
\toprule
\textbf{Name} & \texttt{crypto} \\
\textbf{Accepts} & \texttt{application/json} (exchange API format) \\
\textbf{Input} & Ticker data: \texttt{\{symbol, price, volume,
  timestamp\}} \\
\textbf{Transform} & Entity ID = trading pair (e.g.,
  \texttt{BTC-USDT}). Payload includes price, volume, bid/ask
  spread. Observations batched per window (configurable). \\
\textbf{Streaming} & Yes. Connects to exchange WebSocket feeds.
  Built-in support for: Binance ticker stream, CoinGecko polling,
  generic exchange WebSocket (configurable URL + message format). \\
\textbf{Config} & \texttt{exchange}, \texttt{symbols} (list or
  \texttt{"*"}), \texttt{window\_seconds},
  \texttt{websocket\_url}. \\
\bottomrule
\end{tabularx}

\subsection{A05: HTTP Polling Adapter}

\begin{tabularx}{\textwidth}{l X}
\toprule
\textbf{Name} & \texttt{http\_poll} \\
\textbf{Accepts} & N/A (streaming adapter only) \\
\textbf{Input} & Periodically polls a configurable HTTP endpoint \\
\textbf{Transform} & Response body passed to auto-detected sub-adapter
  (JSON or CSV). \\
\textbf{Streaming} & Yes. Polls at configurable interval. \\
\textbf{Config} & \texttt{url}, \texttt{interval\_seconds},
  \texttt{headers} (optional), \texttt{response\_format}. \\
\bottomrule
\end{tabularx}

\subsection{A06: WebSocket Listener Adapter}

\begin{tabularx}{\textwidth}{l X}
\toprule
\textbf{Name} & \texttt{websocket} \\
\textbf{Accepts} & N/A (streaming adapter only) \\
\textbf{Input} & Connects to an external WebSocket and receives
  messages \\
\textbf{Transform} & Each message parsed as JSON observation. \\
\textbf{Streaming} & Yes. Maintains persistent connection with
  auto-reconnect. \\
\textbf{Config} & \texttt{url}, \texttt{subscribe\_message}
  (optional JSON to send on connect), \texttt{reconnect\_delay\_ms}. \\
\bottomrule
\end{tabularx}

\subsection{A07: File Watcher Adapter}

\begin{tabularx}{\textwidth}{l X}
\toprule
\textbf{Name} & \texttt{file\_watch} \\
\textbf{Accepts} & N/A (streaming adapter only) \\
\textbf{Input} & Watches a directory for new or modified files \\
\textbf{Transform} & New files passed to auto-detected sub-adapter
  (JSON, CSV, syslog). \\
\textbf{Streaming} & Yes. Uses filesystem notifications. \\
\textbf{Config} & \texttt{directory}, \texttt{pattern} (glob),
  \texttt{recursive}. \\
\bottomrule
\end{tabularx}

\subsection{A08: Stdin / Pipe Adapter}

\begin{tabularx}{\textwidth}{l X}
\toprule
\textbf{Name} & \texttt{stdin} \\
\textbf{Accepts} & N/A (streaming adapter only) \\
\textbf{Input} & Reads from stdin (pipe-friendly) \\
\textbf{Transform} & Each line parsed as JSON observation. \\
\textbf{Streaming} & Yes. Reads until EOF. \\
\textbf{Config} & \texttt{format}: \texttt{json} or \texttt{csv}. \\
\bottomrule
\end{tabularx}

\subsection{A09: Synthetic Adapter (existing, enhanced)}

\begin{tabularx}{\textwidth}{l X}
\toprule
\textbf{Name} & \texttt{synthetic} \\
\textbf{Accepts} & N/A (generator, not a real data source) \\
\textbf{Input} & Generates synthetic observation windows for testing \\
\textbf{Config} & \texttt{scenario}: S-Basic, S-Regime, S-Causal,
  S-Break, S-Scale. \\
\bottomrule
\end{tabularx}


\section{Built-In Emission Adapters}\label{sec:ad:emission}

\subsection{E01: Webhook Adapter}

\begin{tabularx}{\textwidth}{l X}
\toprule
\textbf{Name} & \texttt{webhook} \\
\textbf{Output} & POST crystal/alert/metric JSON to a configurable
  URL on each event. \\
\textbf{Config} & \texttt{url}, \texttt{events} (which event types to
  forward), \texttt{headers} (optional), \texttt{retry\_count}. \\
\bottomrule
\end{tabularx}

\subsection{E02: File Writer Adapter}

\begin{tabularx}{\textwidth}{l X}
\toprule
\textbf{Name} & \texttt{file\_writer} \\
\textbf{Output} & Append crystal/alert JSON to JSONL files in a
  configurable directory. \\
\textbf{Config} & \texttt{directory}, \texttt{rotate}: daily/size. \\
\bottomrule
\end{tabularx}

\subsection{E03: Console / Log Adapter}

\begin{tabularx}{\textwidth}{l X}
\toprule
\textbf{Name} & \texttt{console} \\
\textbf{Output} & Print crystal/alert summaries to stdout. \\
\textbf{Config} & \texttt{verbosity}: compact/full. \\
\bottomrule
\end{tabularx}


\section{Adapter Configuration}\label{sec:ad:config}

\begin{lstlisting}[title={Adapter Configuration (in isls config.yaml)}]
adapters:
  ingestion:
    # Crypto market streaming
    - name: crypto
      exchange: binance
      symbols: ["BTCUSDT", "ETHUSDT", "SOLUSDT"]
      window_seconds: 60
      enabled: true

    # File watcher for log ingestion
    - name: file_watch
      directory: /var/log/myapp/
      pattern: "*.json"
      recursive: false
      enabled: true

    # HTTP polling for external API
    - name: http_poll
      url: https://api.example.com/metrics
      interval_seconds: 300
      headers:
        Authorization: "Bearer ..."
      response_format: json
      entity_id_field: sensor_id
      enabled: false

  emission:
    # Webhook for crystal alerts
    - name: webhook
      url: https://hooks.slack.com/services/xxx
      events: [crystal, alert]
      enabled: true

    # Local file backup
    - name: file_writer
      directory: ~/.isls/exports/
      rotate: daily
      enabled: true
\end{lstlisting}


\section{Custom Adapter Development}\label{sec:ad:custom}

Third-party adapters implement the \texttt{IngestionAdapter} or
\texttt{EmissionAdapter} trait and are registered at startup:

\begin{lstlisting}[title={Custom Adapter Example}]
pub struct MyDatabaseAdapter {
    connection_string: String,
}

impl IngestionAdapter for MyDatabaseAdapter {
    fn name(&self) -> &str { "my_database" }

    fn accepts(&self) -> Vec<String> {
        vec!["application/x-mydb".into()]
    }

    fn transform(
        &self,
        input: &[u8],
        config: &AdapterConfig,
    ) -> Result<Vec<Observation>> {
        let rows: Vec<Row> = deserialize(input)?;
        Ok(rows.into_iter().map(|r| Observation {
            entity_id: r.id.clone(),
            timestamp: r.ts,
            payload: serde_json::to_vec(&r)?,
        }).collect())
    }

    fn health(&self) -> AdapterHealth {
        AdapterHealth::Ok
    }
}

// Register at startup:
registry.register_ingestion(Box::new(MyDatabaseAdapter {
    connection_string: "...".into(),
}));
\end{lstlisting}


%%% =====================================================================
\part{Integration}\label{part:integration}

\section{CLI Extension}\label{sec:cli}

\begin{lstlisting}
# Start the gateway server
isls serve --port 8420

# Start with specific adapters enabled
isls serve --port 8420 --adapter crypto --adapter file_watch

# Start engine + gateway in one command
isls serve --port 8420 --run --mode shadow --ticks 0  # 0 = infinite

# Ingest via CLI (uses gateway internally)
isls ingest --adapter json --file data.json
isls ingest --adapter csv --file sensors.csv --entity-id-column device_id
isls ingest --adapter crypto --exchange binance --symbols BTCUSDT,ETHUSDT

# Pipe data into ISLS
cat events.jsonl | isls ingest --adapter stdin --format json
\end{lstlisting}


\section{Typical Deployment Patterns}\label{sec:deploy}

\subsection{Pattern 1: Local Analysis}

\begin{lstlisting}
isls ingest --adapter csv --file market_data.csv
isls run --mode shadow --ticks 500
isls report html
\end{lstlisting}

\subsection{Pattern 2: Live Monitoring}

\begin{lstlisting}
isls serve --port 8420 --run --mode shadow --ticks 0 \
  --adapter crypto --adapter webhook
# Crypto data streams in, crystals stream out to Slack
\end{lstlisting}

\subsection{Pattern 3: Embedded Service}

\begin{lstlisting}
# Another application sends data via REST
curl -X POST http://localhost:8420/ingest \
  -H "Content-Type: application/json" \
  -d '[{"id":"sensor-1","temp":22.5,"ts":1710000000}]'

# Query crystals via REST
curl http://localhost:8420/crystals?limit=5

# Subscribe to live events
wscat -c ws://localhost:8420/events
\end{lstlisting}

\subsection{Pattern 4: Pipeline Integration}

\begin{lstlisting}
# Kafka -> ISLS -> Webhook
kafka-console-consumer --topic events | \
  isls ingest --adapter stdin --format json

# Or use http_poll to pull from any REST API
\end{lstlisting}


%%% =====================================================================
\part{Acceptance Tests}\label{part:tests}

\begin{longtable}{l l p{6.5cm}}
\toprule
\textbf{ID} & \textbf{Name} & \textbf{Procedure} \\
\midrule
\endfirsthead
\midrule
\endhead
\bottomrule
\endfoot
AT-GW1 & Health endpoint & Start gateway; GET /health; verify version
  and uptime. \\
AT-GW2 & JSON ingest & POST /ingest with 100 JSON observations; verify
  receipt shows 100 accepted. \\
AT-GW3 & CSV ingest & POST /ingest/file with CSV; verify observations
  created with correct entity IDs. \\
AT-GW4 & Dedup on ingest & POST same observations twice; verify second
  batch shows duplicates\_skipped > 0. \\
AT-GW5 & Crystal query & Ingest + run; GET /crystals; verify non-empty
  list with correct fields. \\
AT-GW6 & Crystal validate & GET /crystals/\{id\}/validate; verify 8
  checks all pass. \\
AT-GW7 & Manifest query & GET /manifests/\{run\_id\}; verify manifest
  JSON contains expected fields. \\
AT-GW8 & WebSocket events & Connect to /events; run 10 ticks; verify
  at least one \texttt{tick} event and one \texttt{crystal} event
  received. \\
AT-GW9 & Subscription filter & Subscribe to \texttt{["crystal"]} only;
  verify no \texttt{tick} events received. \\
AT-GW10 & Rate limiting & Send observations exceeding rate limit;
  verify 429 response. \\
AT-GW11 & Engine start/stop & POST /engine/start; verify running; POST
  /engine/stop; verify stopped. \\
AT-GW12 & Execute via API & POST /execute with a crystal; verify
  crystals produced. \\
AT-GW13 & Report endpoint & GET /report/html; verify valid HTML
  returned. \\
AT-A1 & JSON adapter & Transform 50 JSON objects; verify 50
  observations with correct entity IDs. \\
AT-A2 & CSV adapter & Transform 100-row CSV; verify 100 observations
  with correct columns. \\
AT-A3 & Syslog adapter & Transform 20 syslog lines; verify 20
  observations with hostname as entity ID. \\
AT-A4 & Crypto adapter & Transform Binance-format ticker JSON; verify
  observations with symbol as entity ID. \\
AT-A5 & Adapter determinism & Transform same input twice; verify
  bit-identical observations. \\
AT-A6 & Adapter registry & Register custom adapter; verify it appears
  in \texttt{list\_ingestion}. \\
AT-A7 & Auto-detect & Submit JSON with Content-Type
  application/json; verify JSON adapter selected. Submit CSV; verify
  CSV adapter selected. \\
AT-A8 & Emission webhook & Configure webhook adapter; emit crystal;
  verify POST received at mock server. \\
AT-A9 & File watcher & Drop a JSON file in watched directory; verify
  observations ingested within 5 seconds. \\
AT-A10 & Stdin pipe & Pipe 10 JSON lines; verify 10 observations
  ingested. \\
\end{longtable}


%%% =====================================================================
\part{Dependency Graph and Build Order}\label{part:deps}

\section{Updated Dependency Graph}

\begin{lstlisting}[title={Complete ISLS Dependency Graph after Phase 4}]
isls-types        <-- (all crates)
    ^
isls-registry
    ^
isls-observe  <-- isls-persist <-- isls-topology (C16)
                       ^                ^
                       |          isls-extract
                       |                ^
                       |          isls-scale (C18)
                       |                ^
               isls-consensus <-- isls-carrier
                       ^
               isls-archive <-- isls-morph
                       ^
               isls-manifest
                       ^
               isls-capsule
                       ^
               isls-scheduler
                       ^
               isls-engine
                       ^
               isls-store (C17)
                       ^
               isls-adapters (C20: adapter trait + 12 built-in adapters)
                       ^
               isls-gateway (C19: REST + WebSocket server)
                       ^
               isls-harness (extended)
                       ^
               isls-cli (extended: serve command)
\end{lstlisting}

\section{External Dependencies}

\begin{tabular}{l l l}
\toprule
\textbf{Crate} & \textbf{Version} & \textbf{Purpose} \\
\midrule
\texttt{axum} & 0.7 & HTTP server framework \\
\texttt{tokio} & 1.0 & Async runtime \\
\texttt{tokio-tungstenite} & 0.24 & WebSocket support \\
\texttt{tower-http} & 0.6 & CORS, compression, timeouts \\
\texttt{notify} & 7.0 & Filesystem watcher (A07) \\
\texttt{csv} & 1.3 & CSV parsing (A02) \\
\texttt{tungstenite} & 0.24 & WebSocket client (A06) \\
\texttt{reqwest} & 0.12 & HTTP client (A05, E01) \\
\bottomrule
\end{tabular}

\section{Build Order}

\begin{enumerate}[nosep]
  \item \texttt{isls-adapters} (C20): depends on \texttt{isls-types},
        \texttt{isls-observe}. New deps: \texttt{csv}, \texttt{notify},
        \texttt{reqwest}, \texttt{tungstenite}.
  \item \texttt{isls-gateway} (C19): depends on \texttt{isls-types},
        \texttt{isls-engine}, \texttt{isls-store},
        \texttt{isls-archive}, \texttt{isls-manifest},
        \texttt{isls-capsule}, \texttt{isls-adapters}. New deps:
        \texttt{axum}, \texttt{tokio}, \texttt{tokio-tungstenite},
        \texttt{tower-http}.
  \item CLI extension: add \texttt{serve} subcommand.
  \item Harness: add AT-GW and AT-A tests.
\end{enumerate}

\section{Test Count}

\begin{tabular}{l r r}
\toprule
\textbf{Phase} & \textbf{New Tests} & \textbf{Cumulative} \\
\midrule
Core (C1--C9) & 130 & 130 \\
Phase 1 (C10--C15) & 31 & 161 \\
Phase 2 (C16--C17) & 20 & 181 \\
Phase 3 (C18) & 15 & 196 \\
Phase 4 (C19--C20) & 23 & $\ge$ 219 \\
\bottomrule
\end{tabular}


%%% =====================================================================
\part{Implementation Handoff}\label{part:handoff}

\section{Verification Sequence}

\begin{lstlisting}
# Build
cargo build --workspace --release 2>&1 | grep -c warning  -> 0

# Test
cargo test --workspace -- --test-threads=1  -> >= 219 passed

# Start gateway
cargo run --bin isls --release -- serve --port 8420 &

# Ingest via REST
curl -s -X POST http://localhost:8420/ingest \
  -H "Content-Type: application/json" \
  -d '[{"id":"e1","value":1.0},{"id":"e2","value":2.0}]'
# -> {"ok":true,"data":{"observations_accepted":2,...}}

# Start engine via REST
curl -s -X POST http://localhost:8420/engine/start \
  -H "Content-Type: application/json" \
  -d '{"mode":"shadow","ticks":50}'

# Query crystals
curl -s http://localhost:8420/crystals | python -m json.tool

# Get HTML report
curl -s http://localhost:8420/report/html > report.html

# WebSocket test
wscat -c ws://localhost:8420/events
\end{lstlisting}

\section{Completion Criteria}

Phase 4 is complete when:
\begin{enumerate}[nosep]
  \item Gateway starts, serves /health, and shuts down cleanly.
  \item All 9 built-in ingestion adapters transform data correctly.
  \item All 3 emission adapters emit events correctly.
  \item REST API returns crystals, manifests, metrics, topology.
  \item WebSocket delivers live events with subscription filtering.
  \item \texttt{isls serve} command works as single entry point.
  \item Custom adapter registration works.
  \item Total tests $\ge 219$, zero warnings.
\end{enumerate}


\vfill
\begin{center}
\rule{0.5\textwidth}{0.4pt}\\[0.5cm]
{\small Generated for \ISLS{} Extension Phase 4 v1.0.0.\\
One gateway. Nine adapters. Three emission sinks.\\
Plug and play.}
\end{center}

\end{document}
